\section{Turning a reading difficulty into a precise question}\label{appE:sec:question}
Suppose you understand what a theorem says but cannot follow one line of its proof. Asking an assistant, or a classmate, to explain the whole proof again rarely helps: the new explanation may skip the same step, or replace it by a different argument that you must then check from the beginning. It works better to turn the difficulty into a precise question, as Habit~4 (Section~\ref{ch04:sec:promises}) suggests. Write down three things:
\begin{itemize}
\item the exact line that you do not follow;
\item what is known at that point of the proof: the hypotheses, and the facts already established;
\item the last step that you do understand.
\end{itemize}
Then ask which definition or earlier fact leads from what is known to the line. A good answer is usually short, one calculation written out in more detail or one small example, and rarely a new technique. Here are two such questions.
\begin{itemize}
\item In Case~1 of the inductive proof of Hall's theorem (Section~\ref{ch09:sec:inductive}), $S$ is a nonempty proper group of students with $|N(S)|=|S|$, and $T$ is any group of the remaining students. A precise question is: ``I know that the task sets $N(S)$ and $N(T)$ can overlap. Why is $|N(S)\cup N(T)|=|N(S)|+|N(T)\setminus N(S)|$?'' The answer is a splitting: every task in the union lies in exactly one of the sets $N(S)$ and $N(T)\setminus N(S)$, and the sizes of two disjoint finite sets add (Section~\ref{ch03:sec:sets}). Nothing more about graphs is needed. Exercise~\ref{ex:9.3} asks you to write this out in full.
\item From the proof of Theorem~\ref{thm:inexact} (Section~\ref{ch18:sec:onestep}): ``I see that each update error is later multiplied by $q$. Why is the weight of $\eps_j$ in the bound for $e_n$ exactly $q^{n-1-j}$?'' The answer is a count. The error $\eps_j$ is made in the update that produces $y_{j+1}$, and each of the $n-(j+1)=n-1-j$ later updates, which produce $y_{j+2},\ldots,y_n$, contributes one factor $q$.
\end{itemize}

Once you have the answer, write the step in your own words. Then change one small feature and work out what changes.

In the Hall example, the proof continues as follows. The students of $S\cup T$ together accept exactly the tasks in $N(S)\cup N(T)$, so Hall's condition gives $|N(S)\cup N(T)|\ge|S\cup T|$, and $|S\cup T|=|S|+|T|$ because the two groups share no student. The identity from the question, with $|N(S)|=|S|$, turns the left side into $|S|+|N(T)\setminus N(S)|$, and subtracting $|S|$ leaves $|N(T)\setminus N(S)|\ge|T|$. Now suppose $T$ may share students with $S$. The identity still holds, but $|S\cup T|=|S|+|T|$ does not, because a student in both groups is counted once on the left and twice on the right. The conclusion no longer follows, and it can even be false: for $T=S$ it would read $0\ge|S|$, which fails because $S$ is nonempty.

In the second example, if the update errors were numbered from $1$, so that $\eps_j$ bounds the error in the update that produces $y_j$, then $n-j$ updates follow it, and its weight becomes $q^{n-j}$. Or suppose $q=1$, so that the map never increases distances but need not shrink them. If it still has a fixed point $x_*$, the same estimate and count work, but every weight is $1$, so the bound is $e_0+\eps_0+\cdots+\eps_{n-1}$. When every update error equals the same $\eta>0$, this is $e_0+n\eta$, which grows without limit, and the error really can grow this fast: every point is a fixed point of the identity map $T(x)=x$ on $\RR$, so take $x_*=0$, start at $y_0=0$, and set $y_{n+1}=y_n+\eta$; then $e_n=n\eta$, exactly the bound. If you can answer variations like these, the explanation has become your own: you can reason with it, not only repeat it.

The same approach works for a definition that you do not understand. Ask, in this order: What kind of object does it apply to? What exactly must such an object satisfy? What is one example, and what is a \emph{near miss}, an object that fails only one of the requirements? For the definition of an odd integer (Definition~\ref{def:parity}), the objects are integers, and $a$ is odd when $a=2k+1$ for some integer $k$. The number $7=2\cdot3+1$ is an example. The number $8$ is a near miss: $8=2\cdot3.5+1$ has the right shape, but $3.5$, the only $k$ with $2k+1=8$, is not an integer, so the one requirement that fails is ``$k$ is an integer''. Finally, when the definition appears in a proof, decide which of the two directions of Section~\ref{ch01:sec:definition} is used: does the proof \emph{unpack} the definition for an object it is given, or \emph{recognize} that a new object satisfies it? The first proof of the book, Proposition~\ref{thm:odd-sum}, does both: it unpacks the fact that $a$ and $b$ are odd, and recognizes that $a+b$ is even.
